\documentclass{article}
\usepackage{amsmath,amssymb,hyperref}
\title{Module extension and the pure-boundary exactness obligation}
\date{15 November 2026}
\begin{document}\maketitle
\paragraph{Status.} Scoped supporting note, not a claimed completion of the relative factory. The standard monoid--bimodule construction below isolates a compatibility question before any passage to a general bicategory or bilimit completion. Its ordinary algebra is background machinery, not a novelty claim.

\paragraph{Independent data.} Let $V,W$ be monoidal categories with reflexive coequalizers preserved separately by tensor, and let $T:V\to W$ be lax monoidal with invertible unit comparison. It sends monoids and their bimodules to monoids and bimodules by applying its tensor comparison before the original action. For composable bimodules $M,N$ over a middle monoid $A$, the composite
\[
 TM\otimes TN\longrightarrow T(M\otimes N)\longrightarrow T(M\otimes_A N)
\]
balances the two $TA$-actions. Indeed the lax monoidal associativity and action naturality identify the two composites from $TM\otimes TA\otimes TN$ after applying the quotient map. Hence it induces a canonical map
\[
 TM\otimes_{TA}TN\longrightarrow T(M\otimes_A N).
\]
Associativity follows by precomposing both candidate maps with the common threefold tensor quotient, an epimorphism; lax monoidal associativity then identifies them. Units use the canonical module isomorphisms $M\otimes_A A\cong M$ and $A\otimes_A N\cong N$. This is a lax extension, not an assertion that the displayed comparison is invertible or that $T$ preserves coequalizers.

\paragraph{Representable exactness is an additional condition.} Consider $V=W=\mathrm{Set}$ with cartesian tensor and $T=\mathcal P_+$, nonempty powerset. Its unit comparison is an isomorphism, and its lax tensor comparison sends $(U,V)$ to $U\times V$. Let $A=1$, let $B$ be a two-element monoid, and let $M$ be a two-element right $1$-set. Regard $B$ as the extension bimodule associated to the monoid map $1\to B$. The preceding comparison becomes
\[
 \mathcal P_+(M)\times\mathcal P_+(B)
 \longrightarrow \mathcal P_+(M\times B).
\]
It has only nonempty rectangles in its image, and misses the two-point diagonal after labeling the two sets. Thus even the extension bimodule of a monoid homomorphism does not automatically give an invertible comparison. Identity-module exactness is weaker than exactness for a chosen pure-interface class.

\paragraph{Exact relevance and limits.} This example does not refute extending a pre-existing pure-exact relative doctrine: powerset already fails product comparison on the displayed pair. It refutes inferring exactness for \emph{newly admitted} pure algebra maps merely from their representability. The identity-monad embedding into a bimodule bicategory retains the old homs and does not create adjoints for arbitrary original pure arrows. One must specify which new pure maps are admitted and prove their base-change comparisons invertible and coherent. The construction also does not supply all required synchronization bilimits.

\paragraph{Even strong tensor preservation is insufficient.} Let $T(Y)=Y^{\mathbb N}$ with its canonical strong Cartesian monoidal structure. Let $A=\{1,a,b\}$ with $a^2=a,b^2=b,ab=b,ba=a$; products of nonidentity elements equal their rightmost factor. Define the left $A$-set $X=\{z_n,x_n:n\in\mathbb N\}$ by fixing every $z_n$, and $a x_n=z_n$, $b x_n=z_{n+1}$. The action equations follow from the multiplication table. For the trivial right $A$-set $1$, the quotient $1\otimes_A X$ is singleton: its generating undirected graph is the connected ray $z_0-x_0-z_1-x_1-\cdots$. But
\[
 1\otimes_{A^{\mathbb N}}X^{\mathbb N}\longrightarrow(1\otimes_A X)^{\mathbb N}=1
\]
is not injective. The sequences $(z_0)_n$ and $(z_n)_n$ lie in different source classes. Each action step or its reverse changes each coordinate by at most one graph edge. A chain of $k$ steps therefore gives coordinate distances bounded by $k$, whereas $d(z_0,z_n)=2n$ is unbounded. This uses the extension bimodule of $A\to1$, not the restriction orientation. All module composites exist; the failure is transport of the quotient, despite strong tensor comparison.

\paragraph{A composition-closed exact pure subcategory.} For the normal lax module extension above, call a monoid map $f:A\to B$ exact if, for every right $A$-module $M$, the canonical comparison
\[
 E_f(M):TM\otimes_{TA}TB\longrightarrow T(M\otimes_A B)
\]
is invertible. These maps form a wide subcategory containing monoid isomorphisms. Identities follow from regular-module units. If $f:A\to B$ and $g:B\to C$ are exact, write the comparison for $gf$ using
\[
 TM\otimes_{TA}TC\cong (TM\otimes_{TA}TB)\otimes_{TB}TC.
\]
Apply $E_f(M)\otimes\mathrm{id}$, then $E_g(M\otimes_A B)$, then $T$ of extension-of-scalars associativity. Lax module associativity identifies this invertible composite with $E_{gf}(M)$. The companion-composition comparison is regular-module contraction, hence invertible; no general tensor-preservation assumption is needed for this closure argument. Isomorphisms follow by transport of actions. This does not claim closure under arbitrary pullbacks, all Morita equivalences, or the existence of synchronization limits. It provides a precise admissible pure class rather than assuming every new algebra map is exact.

\paragraph{Independent local test in the strong case.} If $T$ is strong monoidal, exactness of $f:A\to B$ is equivalent to preservation, for every right $A$-module $M$, of the specific reflexive coequalizer
\[
 M\otimes A\otimes B\rightrightarrows M\otimes B\longrightarrow M\otimes_A B.
\]
The strong tensor identifications identify the first two objects and maps with the pair defining $TM\otimes_{TA}TB$. Its universal comparison is exactly $E_f(M)$, proving both implications. Thus this admission test refers to an independently specified class of local colimits, not to the desired completed crossing or global interaction law. The equivalence is not asserted for merely lax $T$ without further tensor-comparison hypotheses.

\paragraph{Exact established machinery to reuse.} Shulman, \emph{Framed bicategories and monoidal fibrations}, \url{https://arxiv.org/html/0706.1286}, Example 6.2, already constructs the normal lax double functor on bimodules without requiring coequalizer preservation. The calculation above is a specialization, not a new construction. Section 6 distinguishes the restriction comparison, automatically invertible for lax framed functors, from extension comparisons. Theorem 11.5 supplies the general framed module construction under local coequalizer hypotheses. These results replace a proposed reinvention. They do not authorize replacing the original weaker, possibly nonfaithful relative pure interface with a framed equipment having all companions and conjoints.

\paragraph{Source-led narrowing.} Brown's \emph{2-categorical fam constructions}, \url{https://doi.org/10.25949/26794186.v1}, Appendix 9.1, explicitly warns that general modules/comodules need not compose without representing multi-modulations; local completeness hypotheses or suitably free legs are needed. The monoidal calculation here uses explicit coequalizer hypotheses precisely to avoid that omission. Its relation to the original relative biprofunctor factory, the full bicategorical extension and the desired weighted-limit universal property remain proof obligations. A normal lax extension is the correct candidate; upgrading it to a pseudofunctor while retaining locally fully faithful recovery would erase genuinely directed comparison cells.
\end{document}
